\title{Resolvendo problemas - Stewart 9 ed. - Cap. 4}
\author{Túlio Gomes}
\date{may 2026}

\maketitle

\section{Exercísios 4.1}
\subsection{51-66}
Encontre os valores máximo e mínimo das funções no intervalo dado \\
\textbf{51} $f(x) = 12 + 4x -x^2$, $[0, 5]$

\textbf{Solução:} O primeiro passo segundo o método do intervalo fechado é achar os pontos críticos.
\begin{equation}
    \frac{d}{dx}f(x) = -2x + 4
\end{equation}
\begin{equation}
    \begin{split}        
        -2x + 4 = 0 \\
        x = 2
    \end{split}
\end{equation}
O ponto $P=(2, 16)$ é o único ponto crítico.
O segundo passo seria encontra o valor de f nos extremos do intervalo:
\begin{equation}
    \begin{split}
        f(0) = 12 \\
        f(5) = 7
    \end{split}
\end{equation}

Logo, conclui-se que o valor mínimo de $f$ em $[0, 5]$ é 7, e seu máximo é 16.

\subsection{67} 
Se $a$ e $b$ são números positivos, ache o valor máximo de \\
 $f(x) = x^a(1-x)^b$, $0 \leqslant x \leqslant 1$

 \textbf{Solução:} Seguindo o método do intervalo fechado: 
 \begin{equation}
    \begin{split}
        g(x) = x^a \\
        h(x) = (1-x)^b        
    \end{split}
 \end{equation}

 \begin{equation}
    \begin{split}
        g'(x) = \frac{d}{dx}e^{ln(x^a)} \\
        g'(x) = x^a\frac{d}{dx}aln(x) \\
        g'(x) = a\frac{x^a}{x} = ax^{a-1}
    \end{split}
 \end{equation}

  \begin{equation}
    \begin{split}
        h'(x) = \frac{d}{dx}e^{bln(1-x)} \\
        h'(x) = b(1-x)^b\frac{d}{dx}ln(1-x) \\
        h'(x) = -b\frac{(1-x)^b}{1-x} = -b(1-x)^{b-1}
    \end{split}
 \end{equation}

 \begin{equation}
    \begin{split}
        \frac{dy}{dx} = g'(x)h(x) + g(x)h'(x) \\
        =ax^{a-1}(1-x)^b-b(1-x)^{b-1}x^a
    \end{split}
 \end{equation}

 Agora que temos a derivada de $f$, podemos encontrar os pontos críticos.

 \begin{equation}
    \begin{split}
        ax^{a-1}(1-x)^b-b(1-x)^{b-1}x^a = 0 \\
    \end{split}
 \end{equation}

 Isso só é verdade em dois cenários: 
 \begin{equation}
    ax^{a-1}(1-x)^b = bx^a(1-x)^{b-1}
    \label{eq:67first-scenario}
 \end{equation}

 \begin{equation}
    \left\{
        \begin{aligned}
            ax^{a-1}(1-x)^b = 0 \\
            b(1-x)^{b-1}x^a = 0
        \end{aligned}
    \right.
    \label{eq:67second-scenario}
 \end{equation}

 Desenvolvendo a equação~\ref{eq:67first-scenario}:
 \begin{equation}
    \begin{split}
        \frac{a}{b}\frac{x^{a-1}}{x^a} = \frac{(1-x)^{b-1}}{(1-x)^b} \\
        \frac{a}{b}x^{-1} = (1-x)^{-1} \\
        \frac{a}{b} = \frac{x}{(1-x)}
    \end{split}
 \end{equation}

 Então quando $a=x$ e $b=1-x$ $\forall x\in[0, 1]$, esse x será considerado ponto crítico de $f$

 Agora desenvolvendo a primeira equação de~\ref{eq:67second-scenario}:
 \begin{equation}
    \begin{split}
        \left\{
            \begin{aligned}
                ax^{a-1}=0 \\
                ou \\
                (1-x)^b=0
            \end{aligned}
        \right. \\
    \end{split} 
 \end{equation}
 Como $a$ e $b$ são valores positivos, não há solução para esse sistema.
 
 Agora, achando os valores nos extremos do intervalo definido para $x$:
\begin{equation}
    \begin{split}
        f(0) = 0 \\
        f(1) = 0
    \end{split}
\end{equation}
Portanto o valor máximo de $f(x)$ é o maior valor possível de $m$ tal que $m=0$ ou $m=f(c)$, $\forall c\in [0,1]$, tal que $a=c$ e $b=1-c$

\subsection{73}
Após o consumod e bebida alcoólica, a concentração de álcool na corrente sanguínea (CAS) aumenta rapidamente à medida que o álcool é absorvido, seguida por uma diminuição gradual conforme o álcool é metabolizado. A função
\begin{equation}
    C(t) = 0.135te^{-2.802t}
\end{equation}
modela a CAS média em mg/mL, de um grupo de oito indivíduos do seco masculino t horas depois do consumo rápido de 15 mL de etanol puro. Qual é a CAS média máxima nas primeiras 3 horas? Quando ela ocorre?

\textbf{Solução:} já que queremos a média máxima nas primeiras 3 horas, o intervalo da função modeladora deve ser o intervalo fechado $[0, 3]$. Agora devemos utilizar os métodos já conhecidos para determinar o máximo dessa função dado o intervalo.
\begin{equation}
    \begin{split}
        g(t)=0.135t \\
        h(t)=e^{-2.802t}
    \end{split}
\end{equation}
\begin{equation}
    \begin{split}
        g'(t)=0.135 \\
        h'(t)=-2.802e^{-2.802t}
    \end{split}
\end{equation}
\begin{equation}
    \begin{split}
        C'(t) = 0.135t(-2.802e^{-2.802t}) + 0.135e^{-2.802t} \\
        = 0.135e^{-2.802t}(-2.802t + 1)
    \end{split}
\end{equation}
t é ponto crítico de C(t) se C'(t) = 0:
\begin{equation}
    \begin{split}
        0.135e^{-2.802t}(-2.802t + 1) = 0 \\
        -2.802t + 1 = 0 \\
        t = 1/2.802
    \end{split}
\end{equation}
\begin{equation}
    C(1/2.802) = \frac{0.135}{2.802e}
\end{equation}
Achando os extremos:
\begin{equation}
    \begin{split}
        C(0) = 0 \\
        C(3) = 0.135 \cdot 3e^{-2.802 \cdot 3} = 0.405e^{-8.406} 
    \end{split}
\end{equation}
Para determinar qual valor é maior, podemos reorganizar como uma inequação:

\begin{equation}
    \begin{split}
        0.405e^{-8.406} > \frac{0.135}{2.802e} \\
        e^{-8.406} \cdot e > \frac{0.135}{2.802 \cdot 0.405} \\
        e^{-7.406} > \frac{135}{1135.81} \\
        -7.406 > ln(\frac{135}{1135.81})
    \end{split}
\end{equation}
Por aproximação, $ln(\frac{135}{1135.81}) \approx -2.12982655333$. Então a inequação é falsa, uma vez que -7.406 é menor do que o valor aproximado. Isso significa que $C(1/2.802) > C(3)$. Além disso, sabemos que $C(1/2.802) > C(0)$. Logo, um valor de máximo local e absoluto dessa função no intervalo $[0, 3]$ é $C(1/2.802)$

\section{Provas do método do intervalo fechado}
Queremos provar que podemos achar um valor de máximo \textbf{absoluto} de $f$ no intervalo $[a, b]$ utilizando o método do intervalo fechado.

hipótese: se $f: [a,b] \longrightarrow \mathbb{R} $ e $$L = \{x \mid f'(x) = 0 \text{ ou } f'(x) \text{ é indefinida}\}$$
\begin{equation}
        c=max(L \cup \{a, b\})
\end{equation}
então $f(c) > f(x), \forall x \in [a,b]$

\section{4.2 O Teorema do Valor Médio}
\subsection{Teorema de Rolle}
Demonstre que a equação $x^3 + x - 1 = 0$ Possui exatamente uma solução real.
\textbf{Solução:}
Seja $f(x) = x^3 + x - 1$.
Seguindo o Teorema do Valor Intermediário:
\begin{equation}
    \begin{split}
        f(1) = 1 \\
        f(-1) = -3
    \end{split}
\end{equation}

Então a equação possui pelo menos uma raíz real, uma vez que ela é contínua.

se $f$ possui mais de uma raíz real, então $f(a)=0=f(b)$. Deduz-se que $f(x)$ é contínua e derivável em $[a,b]$ e $(a,b)$, respectivamente, uma vez que $f$ é uma função polinomial. De acordo com o Teorema de Rolle, essa função deve possuir algum $c \in [a,b]$, tal que $f'(c)=0$. Contudo, isso não é possível, uma vez que, considerando que $f'(x) = 3x^2 + 1$:
\begin{equation}
    \begin{split}
        3x^2 + 1 \geq 1
    \end{split}
\end{equation}